% Lean: EconLean/Econometrics/Asymptotics/FirstStepGMM.lean, TimeSeries/{WhiteNoise,KalmanRecursion,YuleWalkerPD}.lean
\chapter{Econometrics: feasible two-step GMM, the Kalman recursion, positive definite autocovariances}
\label{chap:econometrics7}

This chapter completes two-step GMM with a first step and a residual-based weight, proves that the Kalman filter's
mean squared errors follow the Riccati recursion and are minimal among linear filters, studies the recursion's
monotonicity and steady state, and shows that nondegenerate innovations make autocovariance matrices positive definite.

\section{Feasible two-step GMM}
\label{sec:feasible-gmm}

Sources: Hansen (1982) Thm~3.2 and Lemma~4.2; White (1980); Hayashi (2000) Prop.~3.4--3.6; Newey and McFadden (1994)
\S4.3. Linear moments $E[z_i(y_i - x_i^\top\theta_0)] = 0$ with $l$ instruments and $k$ regressors. The laws of large
numbers for sample moments and the CLT for $\sqrt n\,\hat g(\theta_0)$ are hypotheses.

\begin{definition}[Sample moments and the residual-based variance estimator]
\label{def:sample-moments}
\lean{Econ.Asymptotics.sampleG, Econ.Asymptotics.sampleZY, Econ.Asymptotics.sampleMoment, Econ.Asymptotics.sHat, Econ.Asymptotics.sCross, Econ.Asymptotics.sQuad, Econ.Asymptotics.expandS}
\leanok
$\hat G = n^{-1}\sum_iz_ix_i^\top$, $n^{-1}\sum_iz_iy_i$, $\hat g(\theta) = n^{-1}\sum_iz_i(y_i - x_i^\top\theta)$ and
$\hat S(\theta) = n^{-1}\sum_i(y_i - x_i^\top\theta)^2z_iz_i^\top$, with the cross and fourth-order moments
$n^{-1}\sum_iu_ix_{ic}z_iz_i^\top$ and $n^{-1}\sum_ix_{ic}x_{ie}z_iz_i^\top$.
\end{definition}

\begin{lemma}[Expansion and consistency of $\hat S$]
\label{lem:shat-consistent}
\lean{Econ.Asymptotics.sHat_eq_expandS, Econ.Asymptotics.expandS_zero, Econ.Asymptotics.continuous_expandS, Econ.Asymptotics.continuous_sHat, Econ.Asymptotics.sHat_consistent, Econ.Asymptotics.sampleZY_eq, Econ.Asymptotics.continuous_sampleG, Econ.Asymptotics.continuous_sampleZY, Econ.Asymptotics.continuous_sampleMoment}
\leanok
\uses{def:sample-moments}
$\hat S(\theta) = \hat S(\theta_0) - 2M_1[\theta - \theta_0] + M_2[\theta - \theta_0,\theta - \theta_0]$, and
$n^{-1}\sum z_iy_i = \hat G\theta_0 + \hat g(\theta_0)$. Hence if $\hat S(\theta_0)\to_pS$, $M_1$ and $M_2$ converge in
probability and $\tilde\theta\to_p\theta_0$, then $\hat S(\tilde\theta)\to_pS$.
\end{lemma}
\begin{proof}
\leanok
Write $y_i - x_i^\top\theta = u_i - x_i^\top d$ with $d = \theta - \theta_0$ and expand the square. The expansion is
jointly continuous in the moments and $d$ and equals $\hat S(\theta_0)$ at $d = 0$; apply the continuous mapping theorem.
\end{proof}

\begin{theorem}[Feasible two-step efficient GMM]
\label{thm:feasible-gmm}
\lean{Econ.Asymptotics.firstStep, Econ.Asymptotics.secondWeight, Econ.Asymptotics.twoStep, Econ.Asymptotics.measurable_firstStep, Econ.Asymptotics.aemeasurable_firstStep, Econ.Asymptotics.two_step_feasible}
\leanok
\uses{lem:shat-consistent, thm:two-step-gmm, thm:gmm-asymptotic}
The first step uses any weight $\hat W_1\to_pW_1$ with $G^\top W_1G$ invertible; the second uses $\hat S(\tilde\theta)^{-1}$.
If $\hat G\to_pG$ of full column rank, $\sqrt n\,\hat g(\theta_0)\to_dL$, $\hat S(\theta_0)\to_pS\succ0$ and the cross and
fourth-order moments converge, then $\sqrt n(\hat\theta - \theta_0)\to_d AL$ with $A = (G^\top S^{-1}G)^{-1}G^\top S^{-1}$
(law $N(0,(G^\top S^{-1}G)^{-1})$ if $L\sim N(0,S)$), and the $J$ statistic has the efficient $J$ limit with mean $l-k$.
\end{theorem}
\begin{proof}
\leanok
$\sqrt n\,\hat g\to_d L$ gives $\hat g\to_p0$, so the first step is consistent (GMM consistency). By
Lemma~\ref{lem:shat-consistent}, $\hat S(\tilde\theta)\to_pS$. Apply Theorem~\ref{thm:two-step-gmm}, after rewriting
$n^{-1}\sum z_iy_i = \hat G\theta_0 + \hat g(\theta_0)$. Measurability follows from continuity of the sample moments and
measurability of the GMM map.
\end{proof}

\begin{theorem}[A data-generating process satisfying the hypotheses]
\label{thm:feasible-gmm-witness}
\lean{Econ.Asymptotics.wEps, Econ.Asymptotics.sum_wEps, Econ.Asymptotics.sum_wEps_sq, Econ.Asymptotics.two_step_feasible_witness}
\leanok
\uses{thm:feasible-gmm}
With $z_i = x_i = 1$, $\theta_0 = 0$ and $y_i = Z/\sqrt n + \varepsilon_i$, where $Z\sim N(0,1)$ and the design errors
$\varepsilon_0 = \sqrt{n/2} = -\varepsilon_1$ (zero otherwise) sum to $0$ with squares summing to $n$, every hypothesis
holds with $S = W_1 = 1$. The example is just-identified, so it does not exercise the weighting; the overidentified case
is covered by Theorem~\ref{thm:two-step-gmm}'s witness.
\end{theorem}
\begin{proof}
\leanok
For $n\ge2$: $\hat G = 1$, $\hat g(0) = Z/\sqrt n$, $\hat S(0) = Z^2/n + 1$, $M_1 = Z/\sqrt n$, $M_2 = 1$. Since
$Z/\sqrt n\to_p0$, all limits follow by continuous mapping.
\end{proof}

\section{The Kalman filter recursion}
\label{sec:kalman-recursion}

Sources: Kalman (1960); Hamilton (1994) \S13.2 and \S13.5; Anderson and Moore (1979) \S3.1 and \S4.4. State
$x_{t+1} = Fx_t + w_t$, observation $y_t = Hx_t + v_t$, filter $\hat x_{t+1} = F(\hat x_t + K_t(y_t - H\hat x_t))$.

\begin{definition}[Riccati map and covariance recursions]
\label{def:riccati}
\lean{Econ.TimeSeries.Kalman.riccati, Econ.TimeSeries.Kalman.covSeq, Econ.TimeSeries.Kalman.gainCov}
\leanok
\uses{def:kalman}
$\mathrm{ric}(P) = F\,J_P(K^*_P)\,F^\top + Q$ with the Joseph form $J$ and Kalman gain $K^*$. The Kalman recursion is
$P_{t+1} = \mathrm{ric}(P_t)$; for an arbitrary gain sequence, $P^K_{t+1} = F\,J_{P^K_t}(K_t)\,F^\top + Q$.
\end{definition}

\begin{lemma}[Properties of the Riccati map]
\label{lem:riccati-props}
\lean{Econ.TimeSeries.Kalman.posSemidef_congr, Econ.TimeSeries.Kalman.joseph_posSemidef, Econ.TimeSeries.Kalman.riccati_eq, Econ.TimeSeries.Kalman.riccati_posSemidef, Econ.TimeSeries.Kalman.joseph_sub_joseph, Econ.TimeSeries.Kalman.riccati_mono, Econ.TimeSeries.Kalman.covSeq_posSemidef, Econ.TimeSeries.Kalman.covSeq_mono_zero, Econ.TimeSeries.Kalman.continuousAt_riccati}
\leanok
\uses{def:riccati, thm:kalman-joseph, thm:kalman-optimal}
$\mathrm{ric}(P) = F(P - PH^\top S^{-1}HP)F^\top + Q$. It preserves positive semidefiniteness, is monotone in the
Loewner order when $R\succ0$, and is continuous wherever $S$ is invertible. From $P_0 = 0$ the recursion is nondecreasing.
\end{lemma}
\begin{proof}
\leanok
The Joseph form is a sum of congruences. For a fixed gain, $J_{P'}(K) - J_P(K) = (I-KH)(P'-P)(I-KH)^\top$; so
$\mathrm{ric}(P') - \mathrm{ric}(P) = F[(J_{P'}(K^*_{P'}) - J_P(K^*_{P'})) + (J_P(K^*_{P'}) - J_P(K^*_P))]F^\top$, and both
brackets are positive semidefinite (the second by optimality of $K^*_P$). Monotonicity from zero follows by induction,
since $\mathrm{ric}(0)\succeq0$.
\end{proof}

\begin{theorem}[Mean squared errors and optimality of the Kalman filter]
\label{thm:kalman-mse}
\lean{Econ.TimeSeries.Kalman.filter_error_succ, Econ.TimeSeries.Kalman.memLp_mulVec_apply, Econ.TimeSeries.Kalman.integral_mulVec_apply_mul, Econ.TimeSeries.Kalman.filter_mse, Econ.TimeSeries.Kalman.covSeq_eq_gainCov, Econ.TimeSeries.Kalman.kalman_filter_mse, Econ.TimeSeries.Kalman.kalman_optimal_gains}
\leanok
\uses{lem:riccati-props}
Suppose the initial error and the noises are square integrable, the noises are white
($E[v_sv_t^\top] = R[s=t]$, $E[w_sw_t^\top] = Q[s=t]$, $E[v_sw_t^\top] = 0$) and uncorrelated with the initial error,
whose second moment is $P_0$. Then a filter with any gains has
$E[(x_t-\hat x_t)(x_t-\hat x_t)^\top] = P^K_t$, and the Kalman filter has $P_t$. Moreover $P^K_t\succeq P_t$ for every
gain sequence and every $t$.
\end{theorem}
\begin{proof}
\leanok
The error satisfies $e_{t+1} = F(I-K_tH)e_t - FK_tv_t + w_t$. By induction, $e_t$ is square integrable and uncorrelated
with $v_s$ and $w_s$ for $s\ge t$. The stacked vector $(e_t,v_t,w_t)$ therefore has block-diagonal second moments, and
the second moment of $e_{t+1}$ is $F\,J_{P_t}(K_t)F^\top + Q$. Optimality follows by induction as in
Lemma~\ref{lem:riccati-props}.
\end{proof}

\begin{theorem}[Steady state]
\label{thm:riccati-steady}
\lean{Econ.TimeSeries.Kalman.single_dotProduct_mulVec_single, Econ.TimeSeries.Kalman.riccati_steady_state}
\leanok
\uses{lem:riccati-props}
From $P_0 = 0$ with $Q\succeq0$ and $R\succ0$: if the recursion is bounded above, it converges to $P_\infty\succeq0$ solving
the algebraic Riccati equation $P_\infty = \mathrm{ric}(P_\infty)$. The bound itself (implied by detectability) is a
hypothesis; uniqueness of the stabilising solution is not formalised.
\end{theorem}
\begin{proof}
\leanok
Each quadratic form $v^\top P_tv$ is nondecreasing and bounded, so it converges; entries converge by polarisation. The
limit is symmetric and positive semidefinite, and continuity of $\mathrm{ric}$ at $P_\infty$ gives the fixed point.
\end{proof}

\begin{definition}[Rademacher white noise]
\label{def:white-noise}
\lean{Econ.TimeSeries.WhiteNoise.coin, Econ.TimeSeries.WhiteNoise.coins, Econ.TimeSeries.WhiteNoise.sgn, Econ.TimeSeries.WhiteNoise.rad, Econ.TimeSeries.WhiteNoise.measurable_rad, Econ.TimeSeries.WhiteNoise.rad_mul_self, Econ.TimeSeries.WhiteNoise.integral_rad, Econ.TimeSeries.WhiteNoise.memLp_rad, Econ.TimeSeries.WhiteNoise.integral_rad_mul}
\leanok
On the infinite product of fair coins, the signs $r_i = \pm1$ are bounded with $E[r_ir_j] = [i = j]$.
\end{definition}

\begin{theorem}[Witnesses for the Kalman results]
\label{thm:kalman-witness}
\lean{Econ.TimeSeries.Kalman.riccati_scalar, Econ.TimeSeries.Kalman.riccati_scalar_witness, Econ.TimeSeries.Kalman.wnV, Econ.TimeSeries.Kalman.wnW, Econ.TimeSeries.Kalman.wnX, Econ.TimeSeries.Kalman.wnY, Econ.TimeSeries.Kalman.wnXh, Econ.TimeSeries.Kalman.kalman_mse_witness}
\leanok
\uses{thm:kalman-mse, thm:riccati-steady, def:white-noise}
For $F = H = Q = R = 1$: $\mathrm{ric}(c) = c/(c+1) + 1$, the recursion from $0$ stays in $[0,2]$, and its limit is the
golden ratio $(1+\sqrt5)/2$. A random walk with Rademacher state and measurement noises satisfies every hypothesis of
Theorem~\ref{thm:kalman-mse}, with mean squared error $1$ at $t = 1$.
\end{theorem}
\begin{proof}
\leanok
Direct computation on $1\times1$ matrices; the fixed point solves $c^2 = c + 1$ with $c\ge0$. Independence of the coins
gives the white-noise moments.
\end{proof}

\section{Positive definite autocovariances}
\label{sec:toeplitz-pd}

Sources: Brockwell and Davis (1991) \S5.1; Hamilton (1994) \S3.4.

\begin{theorem}[Nondegenerate innovations give positive definite Toeplitz matrices]
\label{thm:toeplitz-pd}
\lean{Econ.TimeSeries.YuleWalker.toeplitz_transpose, Econ.TimeSeries.YuleWalker.toeplitz_quadForm, Econ.TimeSeries.YuleWalker.toeplitz_posDef, Econ.TimeSeries.YuleWalker.yule_walker_unique}
\leanok
\uses{def:yule-walker, thm:yule-walker}
If each $x_t$ has an innovation $\varepsilon_t$ orthogonal to the past with $E[\varepsilon_tx_t] = \sigma^2 > 0$, every
Toeplitz matrix $[\gamma(i-j)]$ is positive definite. Hence the Yule--Walker system identifies the AR coefficients without
a rank hypothesis.
\end{theorem}
\begin{proof}
\leanok
$a^\top\Gamma a = E[(\sum_ia_ix_{-i})^2]$. If this vanishes, $\sum_ia_ix_{-i} = 0$ a.s.; pairing with
$\varepsilon_{-i_0}$ at the first nonzero coefficient kills the later lags by orthogonality, leaving $a_{i_0}\sigma^2 = 0$.
\end{proof}

\begin{theorem}[White-noise and MA(1) witnesses]
\label{thm:toeplitz-witness}
\lean{Econ.TimeSeries.YuleWalker.white_noise_witness, Econ.TimeSeries.YuleWalker.ma1_witness}
\leanok
\uses{thm:toeplitz-pd, def:white-noise}
Rademacher white noise is an AR process with $\varphi = 0$ and $\sigma^2 = 1$. The MA(1)
$x_t = \varepsilon_t + \varepsilon_{t-1}/2$ has $\gamma(0) = 5/4$ and $\gamma(\pm1) = 1/2$, so its Toeplitz matrices are
tridiagonal and positive definite.
\end{theorem}
\begin{proof}
\leanok
Expand the products and apply $E[r_ir_j] = [i=j]$.
\end{proof}
